\chapter{Complexity — Loops}
%%% generated by the blueprint pipeline from TCSlib.Complexity.ClassP.BinaryArith.Loops
\section{Overview}

Declarations of the Lean module \texttt{TCSlib.Complexity.ClassP.BinaryArith.Loops}.

\section{Declarations}

\begin{theorem}[Composition preserves polynomial time]
\lean{Complexity.BinArith.polyTimeComputable_app}
\label{Complexity.BinArith.polyTimeComputable_app}
\leanok
\uses{Complexity.PolyTimeComputable, Complexity.PolyTimeComputable.comp}
\difficulty{1}
If $f$ and $g$ are polynomial-time computable functions on bit strings, then their
composition $x\mapsto f(g(x))$ is polynomial-time computable.
\end{theorem}

\begin{theorem}[Polynomial-time composition with a bit output]
\lean{Complexity.BinArith.polyTimeComputable_test_app}
\label{Complexity.BinArith.polyTimeComputable_test_app}
\leanok
\uses{Complexity.PolyTimeComputable, Complexity.PolyTimeComputable.comp}
\difficulty{1}
Let $p$ map bit strings to bits and let $g$ map bit strings to bit strings. If the
function $x\mapsto [p(x)]$ is polynomial-time computable and $g$ is polynomial-time
computable, then the function $x\mapsto [p(g(x))]$ is polynomial-time computable.
\end{theorem}

\begin{definition}[Bounded-list iteration step]
\lean{Complexity.BinArith.forStep}
\label{Complexity.BinArith.forStep}
\leanok
\uses{Complexity.isNilB, Complexity.pairFstD, Complexity.pairSndD, Turing.pairEncode}
Given a function $f$ from Boolean lists to Boolean lists and a Boolean list $S$, return
$S$ if its first decoded component is nil. Otherwise, pair the tail of that first
decoded component with the result of applying $f$ to the second decoded component of
$S$.
\end{definition}

\begin{definition}[Bounded-list iteration]
\lean{Complexity.BinArith.forLoop}
\label{Complexity.BinArith.forLoop}
\leanok
\uses{Complexity.BinArith.forStep}
Given a function $f$ from Boolean lists to Boolean lists and a Boolean list $S$, apply
the bounded-list iteration step associated with $f$ exactly $|S|+1$ times to $S$, where
$|S|$ is its length.
\end{definition}

\begin{theorem}[Iteration of the encoded loop step]
\lean{Complexity.BinArith.iterate_forStep}
\label{Complexity.BinArith.iterate_forStep}
\leanok
\uses{Complexity.BinArith.forStep, Complexity.isNilB, Complexity.pairFstD_pairEncode, Complexity.pairSndD_pairEncode, Turing.pairEncode}
\difficulty{5}
Let $f$ be a function from bit strings to bit strings, and let $c,d$ be bit strings.
After applying the loop step $i$ times to the pair encoding of $c$ and $d$, the result
is the pair encoding of the suffix of $c$ after deleting its first $i$ bits and the
result of applying $f$ repeatedly to $d$ exactly $mathrm{min}(i,|c|)$ times.
\end{theorem}

\begin{theorem}[For-loop iteration on an encoded pair]
\lean{Complexity.BinArith.forLoop_pairEncode}
\label{Complexity.BinArith.forLoop_pairEncode}
\leanok
\uses{Complexity.BinArith.forLoop, Complexity.BinArith.iterate_forStep, Turing.length_pairEncode, Turing.pairEncode}
\difficulty{3}
Let $f$ map finite bit strings to finite bit strings, and let $c,d$ be finite bit
strings. Starting from the self-delimiting pair encoding $\langle c,d\rangle$, iterate
the step that leaves the state unchanged when its total first projection is empty, and
otherwise replaces it by the encoding of that projection with its first bit removed and
the result of applying $f$ to its total second projection. After $|\langle
c,d\rangle|+1$ steps, the state is the encoding $\langle\varepsilon,f^{|c|}(d)\rangle$,
where $\varepsilon$ is the empty bit string.
\end{theorem}

\begin{theorem}[For-step on an empty first projection]
\lean{Complexity.BinArith.forStep_of_fst_nil}
\label{Complexity.BinArith.forStep_of_fst_nil}
\leanok
\uses{Complexity.BinArith.forStep, Complexity.isNilB, Complexity.pairFstD}
\difficulty{1}
Let $f$ map finite bit strings to finite bit strings, and let $S$ be a finite bit string
whose total first projection is empty. Then one application of the for-step
transformation leaves $S$ unchanged.
\end{theorem}

\begin{theorem}[Polynomial-Time Computability of a Bounded Iteration]
\lean{Complexity.BinArith.polyTimeComputable_forLoop}
\label{Complexity.BinArith.polyTimeComputable_forLoop}
\leanok
\uses{Complexity.BinArith.forLoop, Complexity.BinArith.forStep, Complexity.BinArith.forStep_of_fst_nil, Complexity.BinArith.iterate_forStep, Complexity.BinArith.polyTimeComputable_iterate, Complexity.PolyTimeComputable, Complexity.PolyTimeComputable.comp, Complexity.PolyTimeComputable.pairEncode, Complexity.pairFstD, Complexity.polyTimeComputable_id, Complexity.polyTimeComputable_isNil, Complexity.polyTimeComputable_ite, Complexity.polyTimeComputable_pairFstD, Complexity.polyTimeComputable_pairSndD, Complexity.polyTimeComputable_tail, Turing.eq_pairEncode_of_pairDecode, Turing.length_pairEncode, Turing.pairDecode, Turing.pairEncode}
\difficulty{6}
Let $f$ be a polynomial-time computable function on bit strings. Suppose there are
natural numbers $B,b$ such that, for all bit strings $c,d$ and all natural numbers
$j\leq |c|$, the $j$-fold iterate of $f$ applied to $d$ has length at most
$B\bigl(|\langle c,d\rangle|+1\bigr)^b$, where $\langle c,d\rangle$ is the
self-delimiting pair encoding. Define a state update that leaves a state unchanged when
its first decoded component is empty, and otherwise removes the first bit of that
component and replaces the second component by its image under $f$. Iterating this
update $|S|+1$ times from $S$ gives a polynomial-time computable function of $S$.
\end{theorem}

\begin{theorem}[Length bound for iterated context-preserving maps]
\lean{Complexity.BinArith.length_iterate_le_of_ctx}
\label{Complexity.BinArith.length_iterate_le_of_ctx}
\leanok
\uses{Complexity.pairFstD}
\difficulty{5}
Let $f$ map finite bit strings to finite bit strings, and let $A$ be a natural number.
Suppose that for every bit string $d$, the first decoded component of $f(d)$ equals the
first decoded component of $d$, and $|f(d)|≤|d|+A(|\pi_1(d)|+1)$, where $\pi_1$ denotes
the first decoded component, taken to be empty if parsing fails. Then for every bit
string $d$ and natural number $j$, $|f^{j}(d)|≤|d|+jA(|\pi_1(d)|+1)$, where $f^j$
denotes $j$-fold iteration.
\end{theorem}

\begin{theorem}[Bound on iterated context length]
\lean{Complexity.BinArith.forLoop_bound_of_ctx}
\label{Complexity.BinArith.forLoop_bound_of_ctx}
\leanok
\uses{Complexity.BinArith.length_iterate_le_of_ctx, Complexity.length_pairFstD_le, Complexity.pairFstD, Turing.length_pairEncode, Turing.pairEncode}
\difficulty{4}
Let $f$ map finite bit strings to finite bit strings, and let $A$ be a natural number.
Suppose that for every bit string $d$, the total first projection of $f(d)$ equals that
of $d$, and $|f(d)|\le |d|+A\bigl(|\pi_1(d)|+1\bigr)$, where $\pi_1$ is the total first
projection for the self-delimiting pair encoding. Then for all bit strings $c,d$ and all
natural numbers $j\le |c|$, the $j$-fold iterate of $f$ at $d$ has length at most
$(A+1)(|\langle c,d\rangle|+1)^2$, where $\langle c,d\rangle$ is the self-delimiting
pair encoding.
\end{theorem}

\begin{definition}[Boolean head transition]
\lean{Complexity.BinArith.headδ}
\leanok
For Booleans $q$ and $b$, the head transition is their conjunction: it is true exactly
when both are true.
\end{definition}

\begin{definition}[Optional Boolean head output]
\lean{Complexity.BinArith.headO}
\label{Complexity.BinArith.headO}
\leanok
For Booleans $q$ and $b$, return the optional value containing true when both are true,
and return no value otherwise.
\end{definition}

\begin{definition}[Unary head transduction]
\lean{Complexity.BinArith.uHead}
\label{Complexity.BinArith.uHead}
\leanok
\uses{Complexity.BinArith.headO, Complexity.BinArith.headδ, Complexity.transduce}
For a Boolean list $w$, let $u(w)$ be the output of the transducer specified by `headδ`
and `headO`, started in state `true` and run on $w$.
\end{definition}

\begin{definition}[Boolean remainder predicate]
\lean{Complexity.BinArith.restδ}
\leanok
Define a function on pairs of Boolean values by declaring its value to be true exactly
when the first value is true or the second value is false.
\end{definition}

\begin{definition}[Rest output rule]
\lean{Complexity.BinArith.restO}
\label{Complexity.BinArith.restO}
\leanok
For Boolean values $q$ and $b$, define the output to be $b$ when $q$ is true, and no
output when $q$ is false.
\end{definition}

\begin{definition}[Unary rest transduction]
\lean{Complexity.BinArith.uRest}
\label{Complexity.BinArith.uRest}
\leanok
\uses{Complexity.BinArith.restO, Complexity.BinArith.restδ, Complexity.transduce}
For a Boolean list $w$, let $u(w)$ be the output of the transducer specified by `restδ`
and `restO`, started in state `false` and run on $w$.
\end{definition}

\begin{definition}[Presence of a false entry]
\lean{Complexity.BinArith.hasFalse}
\label{Complexity.BinArith.hasFalse}
\leanok
For a Boolean list $w$, define a Boolean that is true exactly when $w$ contains at least
one false entry.
\end{definition}

\begin{theorem}[Output from the initial false state]
\lean{Complexity.BinArith.transduce_head_false}
\label{Complexity.BinArith.transduce_head_false}
\leanok
\uses{Complexity.BinArith.headO, Complexity.BinArith.headδ, Complexity.transduce}
\difficulty{4}
For every finite Boolean word $w$, the one-pass transducer whose state is a bit, whose
transition sends $(q,b)$ to $q \wedge b$, and which emits a true bit exactly when $q
\wedge b$ is true produces the empty word on input $w$ when started in the false state.
\end{theorem}

\begin{theorem}[Output from the initial true state]
\lean{Complexity.BinArith.transduce_rest_true}
\label{Complexity.BinArith.transduce_rest_true}
\leanok
\uses{Complexity.BinArith.restO, Complexity.BinArith.restδ, Complexity.transduce}
\difficulty{4}
For every finite Boolean word $w$, the one-pass transducer whose state is a bit, whose
transition sends $(q,b)$ to $q \vee \neg b$, and which emits the input bit $b$ exactly
when the state is true produces $w$ when started in the true state.
\end{theorem}

\begin{theorem}[Output on a false initial bit]
\lean{Complexity.BinArith.uHead_false_cons}
\label{Complexity.BinArith.uHead_false_cons}
\leanok
\uses{Complexity.BinArith.headO, Complexity.BinArith.headδ, Complexity.BinArith.transduce_head_false, Complexity.BinArith.uHead, Complexity.transduce}
\difficulty{1}
For every finite bit string $w$, the output of the unary-head transducer on the string
beginning with $0$ is empty: $uHead(0w)=\varepsilon$.
\end{theorem}

\begin{theorem}[Output on a true initial bit]
\lean{Complexity.BinArith.uHead_true_cons}
\label{Complexity.BinArith.uHead_true_cons}
\leanok
\uses{Complexity.BinArith.headO, Complexity.BinArith.headδ, Complexity.BinArith.uHead, Complexity.transduce}
\difficulty{1}
For every finite bit string $w$, the unary-head transducer outputs an initial $1$
followed by its output on $w$: $uHead(1w)=1uHead(w)$.
\end{theorem}

\begin{theorem}[Output after an initial false bit]
\lean{Complexity.BinArith.uRest_false_cons}
\label{Complexity.BinArith.uRest_false_cons}
\leanok
\uses{Complexity.BinArith.restO, Complexity.BinArith.restδ, Complexity.BinArith.transduce_rest_true, Complexity.BinArith.uRest, Complexity.transduce}
\difficulty{1}
For every finite bit string $w$, the unary-rest transducer maps the string beginning
with $0$ followed by $w$ to $w$: $uRest(0w)=w$.
\end{theorem}

\begin{theorem}[Output after an initial true bit]
\lean{Complexity.BinArith.uRest_true_cons}
\label{Complexity.BinArith.uRest_true_cons}
\leanok
\uses{Complexity.BinArith.restO, Complexity.BinArith.restδ, Complexity.BinArith.uRest, Complexity.transduce}
\difficulty{1}
For every finite bit string $w$, the unary-rest transducer has the same output on $1w$
as on $w$: $uRest(1w)=uRest(w)$.
\end{theorem}

\begin{theorem}[Output on the empty string]
\lean{Complexity.BinArith.uHead_nil}
\label{Complexity.BinArith.uHead_nil}
\leanok
\uses{Complexity.BinArith.uHead}
\difficulty{1}
The unary-head transducer outputs the empty string on the empty input:
$uHead(\varepsilon)=\varepsilon$.
\end{theorem}

\begin{theorem}[Output on the empty string]
\lean{Complexity.BinArith.uRest_nil}
\label{Complexity.BinArith.uRest_nil}
\leanok
\uses{Complexity.BinArith.uRest}
\difficulty{1}
The unary-rest transducer outputs the empty string on the empty input:
$uRest(\varepsilon)=\varepsilon$.
\end{theorem}

\begin{theorem}[False Bits in a Prepending List]
\lean{Complexity.BinArith.hasFalse_cons}
\label{Complexity.BinArith.hasFalse_cons}
\leanok
\uses{Complexity.BinArith.hasFalse}
\difficulty{1}
For a Boolean $b$ and a list $w$ of Booleans, the list formed by prepending $b$ has a
false entry exactly when $b$ is false or $w$ has a false entry.
\end{theorem}

\begin{theorem}[False-bit test on the empty list]
\lean{Complexity.BinArith.hasFalse_nil}
\label{Complexity.BinArith.hasFalse_nil}
\leanok
\uses{Complexity.BinArith.hasFalse}
\difficulty{1}
For a list of Boolean values, define the false-bit test to return true exactly when the
list contains a false entry. On the empty list, this test returns false.
\end{theorem}

\begin{theorem}[Characterization of natural-number decoding]
\lean{Complexity.BinArith.decodeNat_eq}
\label{Complexity.BinArith.decodeNat_eq}
\leanok
\uses{BoolCircuit.decodeNat, Complexity.BinArith.hasFalse, Complexity.BinArith.hasFalse_cons, Complexity.BinArith.uHead, Complexity.BinArith.uHead_true_cons, Complexity.BinArith.uRest, Complexity.BinArith.uRest_true_cons}
\difficulty{4}
For every bit string $w$, decoding it as a natural number and residual word gives the
pair consisting of the length of the output of the head transducer and the output of the
rest transducer, if $w$ contains at least one false bit; otherwise decoding is
undefined.
\end{theorem}

\begin{theorem}[Output on a unary natural encoding]
\lean{Complexity.BinArith.uHead_encodeNat}
\label{Complexity.BinArith.uHead_encodeNat}
\leanok
\uses{BoolCircuit.encodeNat, Complexity.BinArith.uHead}
\difficulty{4}
For every natural number $j$ and finite Boolean word $w$, feeding the one-pass
transducer that starts in the true state, has transition $(q,b)\mapsto q\wedge b$, and
emits true exactly when $q\wedge b$ is true the concatenation of $j$ true bits, one
false bit, and $w$ produces exactly $j$ true bits.
\end{theorem}

\begin{theorem}[Removing a unary natural encoding]
\lean{Complexity.BinArith.uRest_encodeNat}
\label{Complexity.BinArith.uRest_encodeNat}
\leanok
\uses{BoolCircuit.encodeNat, Complexity.BinArith.uRest}
\difficulty{4}
For every natural number $j$ and finite bit string $w$, the unary-rest transducer maps
the unary encoding of $j$ followed by $w$ to $w$: $uRest(1^j0w)=w$.
\end{theorem}

\begin{theorem}[False bit in a natural number encoding]
\lean{Complexity.BinArith.hasFalse_encodeNat}
\label{Complexity.BinArith.hasFalse_encodeNat}
\leanok
\uses{BoolCircuit.encodeNat, Complexity.BinArith.hasFalse}
\difficulty{4}
For every natural number $j$ and bit string $w$, the concatenation of $w$ with the
encoding consisting of $j$ true bits followed by a false bit contains a false bit.
\end{theorem}

\begin{theorem}[Decomposition of a bit string containing zero]
\lean{Complexity.BinArith.eq_encodeNat_of_hasFalse}
\label{Complexity.BinArith.eq_encodeNat_of_hasFalse}
\leanok
\uses{BoolCircuit.encodeNat, Complexity.BinArith.hasFalse, Complexity.BinArith.hasFalse_cons, Complexity.BinArith.uHead, Complexity.BinArith.uHead_true_cons, Complexity.BinArith.uRest, Complexity.BinArith.uRest_true_cons}
\difficulty{5}
Let $w$ be a bit string containing at least one $0$. Apply the first transducer, started
in its specified initial state, to obtain an output string $H$, and apply the second
transducer, also from its specified initial state, to obtain an output string $R$. Then
$w$ is a string of $|H|$ consecutive $1$s followed by $0$ and then $R$.
\end{theorem}

\begin{theorem}[Length bound for the output]
\lean{Complexity.BinArith.length_uRest_le}
\label{Complexity.BinArith.length_uRest_le}
\leanok
\uses{Complexity.BinArith.uRest, Complexity.BinArith.uRest_true_cons}
\difficulty{4}
For every finite bit string $w$, process its bits from left to right with a Boolean
state initially false: on reading $b$, emit $b$ exactly when the current state is true,
and update the state to true if it was already true or if $b=0$. The resulting output
string has length at most $|w|$.
\end{theorem}

\begin{theorem}[Length bound for the unary head output]
\lean{Complexity.BinArith.length_uHead_le}
\label{Complexity.BinArith.length_uHead_le}
\leanok
\uses{Complexity.BinArith.uHead, Complexity.BinArith.uHead_true_cons}
\difficulty{4}
For every bit string $w$, the unary head transducer's output on $w$ has length at most
$|w|$.
\end{theorem}

\begin{theorem}[Polynomial-time unary head transduction]
\lean{Complexity.BinArith.polyTimeComputable_uHead}
\label{Complexity.BinArith.polyTimeComputable_uHead}
\leanok
\uses{Complexity.BinArith.headO, Complexity.BinArith.headδ, Complexity.BinArith.uHead, Complexity.PolyTimeComputable, Complexity.polyTimeComputable_transduce}
\difficulty{1}
The bit-string function obtained by processing the input from left to right with a
Boolean state initially true, updating the state to its conjunction with each input bit,
and emitting a $1$ exactly when both the current state and input bit are true is
polynomial-time computable.
\end{theorem}

\begin{theorem}[Polynomial-time unary rest transduction]
\lean{Complexity.BinArith.polyTimeComputable_uRest}
\label{Complexity.BinArith.polyTimeComputable_uRest}
\leanok
\uses{Complexity.BinArith.restO, Complexity.BinArith.restδ, Complexity.BinArith.uRest, Complexity.PolyTimeComputable, Complexity.polyTimeComputable_transduce}
\difficulty{1}
The bit-string function obtained by processing the input from left to right with a
Boolean state initially false, updating the state to its disjunction with the negation
of each input bit, and emitting that input bit exactly when the current state is true is
polynomial-time computable.
\end{theorem}

\begin{theorem}[Polynomial-Time Detection of a False Bit]
\lean{Complexity.BinArith.polyTimeComputable_hasFalse}
\label{Complexity.BinArith.polyTimeComputable_hasFalse}
\leanok
\uses{Complexity.BinArith.hasFalse, Complexity.BinArith.polyTimeComputable_any, Complexity.PolyTimeComputable, Complexity.PolyTimeComputable.comp, Complexity.polyTimeComputable_transduce, Complexity.transduce}
\difficulty{5}
The function sending a finite Boolean word $w$ to the one-bit word $[\mathsf{true}]$ if
some entry of $w$ is false, and to $[\mathsf{false}]$ otherwise, is polynomial-time
computable.
\end{theorem}

\begin{definition}[First projection after iteration]
\lean{Complexity.BinArith.getW}
\label{Complexity.BinArith.getW}
\leanok
\uses{Complexity.BinArith.forLoop, Complexity.pairFstD, Complexity.pairSndD}
For a list of Booleans $z$, define $w$ to be the first component of the second component
obtained by iterating the second-component projection on $z$ and then taking the second
component once more.
\end{definition}

\begin{theorem}[Projection iteration on encoded words]
\lean{Complexity.BinArith.iterate_pairSndD_encWords}
\label{Complexity.BinArith.iterate_pairSndD_encWords}
\leanok
\uses{Complexity.BinArith.encWords, Complexity.BinArith.iterate_pairSndD_nil, Complexity.pairSndD}
\difficulty{4}
For any finite list $L$ of bit strings and natural number $j$, applying the total second
projection $j$ times to the concatenation of the encodings $\langle w,[]\rangle$ for $w$
in $L$ gives the corresponding concatenation for the list obtained by removing the first
$j$ entries of $L$.
\end{theorem}

\begin{theorem}[Retrieving a Word from an Encoding]
\lean{Complexity.BinArith.getW_pairEncode}
\label{Complexity.BinArith.getW_pairEncode}
\leanok
\uses{Complexity.BinArith.encWords, Complexity.BinArith.forLoop_pairEncode, Complexity.BinArith.getW, Complexity.BinArith.iterate_pairSndD_encWords, Complexity.pairSndD_pairEncode, Turing.pairEncode}
\difficulty{4}
Let $J$ be a bit string and let $L$ be a list of bit strings. Encode $L$ by
concatenating the pair encodings of its words with empty second components. Applying the
word-retrieval procedure to the pair encoding of $J$ and this encoding of $L$ returns
the word at index $|J|$ in $L$, or the empty word if that index is out of range.
\end{theorem}

\begin{theorem}[Retrieving a Natural Number's Binary Representation]
\lean{Complexity.BinArith.getW_pairEncode_bits}
\label{Complexity.BinArith.getW_pairEncode_bits}
\leanok
\uses{Complexity.BinArith.encWords, Complexity.BinArith.getW, Complexity.BinArith.getW_pairEncode, Turing.pairEncode}
\difficulty{4}
Let $J$ be a bit string and let $v_1,\ldots,v_m$ be a finite list of natural numbers.
Encode the list of their binary representations by concatenating the pair encodings of
those representations with empty second components. Applying the word-retrieval
procedure to the pair encoding of $J$ and this encoded list returns the binary
representation of the entry at index $|J|$, taking that entry to be $0$ when the index
is out of range.
\end{theorem}

\begin{theorem}[Polynomial-Time Computability of a Loop Output]
\lean{Complexity.BinArith.polyTimeComputable_getW}
\label{Complexity.BinArith.polyTimeComputable_getW}
\leanok
\uses{Complexity.BinArith.forLoop, Complexity.BinArith.getW, Complexity.BinArith.length_iterate_pairSndD_le, Complexity.BinArith.polyTimeComputable_forLoop, Complexity.PolyTimeComputable, Complexity.PolyTimeComputable.comp, Complexity.pairSndD, Complexity.polyTimeComputable_pairFstD, Complexity.polyTimeComputable_pairSndD, Turing.length_pairEncode}
\difficulty{4}
For a bit string $z$, repeatedly apply the following update $|z|+1$ times: if the first
decoded component is empty, leave the current string unchanged; otherwise remove its
first bit and replace its second decoded component by that component's image under the
total second projection of the pair parser. Then take the first decoded component of the
second decoded component of the resulting string, using the empty string when decoding
fails. The function so defined is computable in polynomial time by a finite Turing
machine.
\end{theorem}

\begin{theorem}[Pair-decoding cases]
\lean{Complexity.BinArith.pair_cases}
\label{Complexity.BinArith.pair_cases}
\leanok
\uses{Complexity.pairFstD, Complexity.pairSndD, Turing.eq_pairEncode_of_pairDecode, Turing.pairDecode, Turing.pairEncode}
\difficulty{4}
For every bit string $z$, either both total projections of $z$ are empty, or $z$ is the
pair encoding of its total first and second projections.
\end{theorem}

\begin{definition}[Length check and equality test]
\lean{Complexity.BinArith.eqW}
\label{Complexity.BinArith.eqW}
\leanok
\uses{Complexity.BinArith.beq, Complexity.pairFstD, Complexity.pairSndD}
For a list $z$, define a Boolean equal to the conjunction of the test that the lengths
of the first and second projections of $z$ are equal, and the test `beq z`.
\end{definition}

\begin{theorem}[Equality test on paired bit strings]
\lean{Complexity.BinArith.eqW_pairEncode}
\label{Complexity.BinArith.eqW_pairEncode}
\leanok
\uses{Complexity.BinArith.beq, Complexity.BinArith.eqW, Complexity.LogProg.eq_of_bitsVal_eq, Complexity.pairFstD_pairEncode, Complexity.pairSndD_pairEncode, Turing.pairEncode}
\difficulty{5}
For any bit strings $a$ and $b$, the equality test applied to their pair encoding
returns true exactly when $a=b$.
\end{theorem}

\begin{theorem}[Polynomial-Time Computability of a Binary Equality Test]
\lean{Complexity.BinArith.polyTimeComputable_eqW}
\label{Complexity.BinArith.polyTimeComputable_eqW}
\leanok
\uses{Complexity.BinArith.eqW, Complexity.BinArith.polyTimeComputable_beq, Complexity.PolyTimeComputable, Complexity.polyTimeComputable_and, Complexity.polyTimeComputable_lenEq}
\difficulty{1}
For a bit string $z$, parse its two components, using the empty string for either
component when parsing fails. The function that outputs the single bit $1$ exactly when
the two components have equal lengths and represent equal natural numbers in
least-significant-bit-first binary, and outputs $0$ otherwise, is computable in
polynomial time by a finite Turing machine.
\end{theorem}

\begin{definition}[Unary multiplication step]
\lean{Complexity.BinArith.mulUStep}
\label{Complexity.BinArith.mulUStep}
\leanok
\uses{Complexity.pairFstD, Complexity.pairSndD, Turing.pairEncode}
For a Boolean list $d$, define a paired encoding whose first component is the first
projection of $d$ and whose second component is the concatenation of the second
projection of $d$ with its first projection.
\end{definition}

\begin{definition}[Unsigned multiplication routine]
\lean{Complexity.BinArith.mulU}
\label{Complexity.BinArith.mulU}
\leanok
\uses{Complexity.BinArith.forLoop, Complexity.BinArith.mulUStep, Complexity.pairFstD, Complexity.pairSndD, Turing.pairEncode}
For a Boolean list $z$, apply the multiplication loop, initialized with the first
component of $z$ and an encoding of its second component paired with the empty list, and
return the second component of the second component of the resulting state.
\end{definition}

\begin{theorem}[Iteration of the multiplication step]
\lean{Complexity.BinArith.iterate_mulUStep}
\label{Complexity.BinArith.iterate_mulUStep}
\leanok
\uses{Complexity.BinArith.mulUStep, Turing.pairEncode}
\difficulty{4}
For any bit strings $b,\mathrm{acc}$ and natural number $j$, applying the multiplication
step $j$ times to $\langle b,\mathrm{acc}\rangle$ gives the pair encoding of $b$ and the
concatenation of $\mathrm{acc}$ with $j$ copies of $b$.
\end{theorem}

\begin{theorem}[Length of the unary product output]
\lean{Complexity.BinArith.length_mulU}
\label{Complexity.BinArith.length_mulU}
\leanok
\uses{Complexity.BinArith.forLoop_pairEncode, Complexity.BinArith.iterate_mulUStep, Complexity.BinArith.mulU, Turing.pairEncode}
\difficulty{1}
For any two finite bit strings $a$ and $b$, the output of the unary multiplication
procedure on their self-delimiting pair encoding has length $|a|\,|b|$, where $|a|$ and
$|b|$ denote the numbers of bits in the strings.
\end{theorem}

\begin{theorem}[Polynomial-Time Computability of Unary Multiplication]
\lean{Complexity.BinArith.polyTimeComputable_mulU}
\label{Complexity.BinArith.polyTimeComputable_mulU}
\leanok
\uses{Complexity.BinArith.forLoop_bound_of_ctx, Complexity.BinArith.mulU, Complexity.BinArith.mulUStep, Complexity.BinArith.pair_cases, Complexity.BinArith.polyTimeComputable_forLoop, Complexity.BinArith.polyTimeComputable_slot3, Complexity.PolyTimeComputable, Complexity.PolyTimeComputable.append, Complexity.PolyTimeComputable.comp, Complexity.PolyTimeComputable.pairEncode, Complexity.polyTimeComputable_const, Complexity.polyTimeComputable_pairFstD, Complexity.polyTimeComputable_pairSndD, Turing.length_pairEncode, Turing.pairEncode}
\difficulty{5}
Encode pairs by doubling each bit of the first component, appending $01$, and then the
second component, and decode inputs by reading aligned pairs, with empty strings as
defaults when decoding fails. Starting from the encoding of the first decoded component
of the input paired with the encoding of its second decoded component paired with the
empty string, iterate the update that preserves the first component and appends it to
the second component exactly $|z|+1$ times; then take the second component twice from
the resulting nested pair. The resulting function is polynomial-time computable.
\end{theorem}

\begin{definition}[Chunking step]
\lean{Complexity.BinArith.chunkStep}
\label{Complexity.BinArith.chunkStep}
\leanok
\uses{Complexity.PrefixByLength.drop, Complexity.PrefixByLength.take, Complexity.isNilB, Complexity.pairFstD, Complexity.pairSndD, Turing.pairEncode}
For a Boolean-list state $d$, return $d$ unchanged if the first component of the second
component of $d$ is empty. Otherwise, keep the first component of $d$ as the first
component of the new state; pair the remainder of the prefix of the second component
with the second component of the second component of $d$, followed by the prefix itself
paired with the empty list, as the second component of the new state.
\end{definition}

\begin{definition}[Chunking routine]
\lean{Complexity.BinArith.chunkW}
\label{Complexity.BinArith.chunkW}
\leanok
\uses{Complexity.BinArith.chunkStep, Complexity.BinArith.forLoop, Complexity.pairFstD, Complexity.pairSndD, Turing.pairEncode}
For a Boolean list $z$, initialize the chunking loop with the second component of $z$,
paired with the first component of $z$, paired with the second component of $z$ and the
empty list. Return the second component of the second component of the second component
of the resulting state.
\end{definition}

\begin{definition}[List of fixed-length chunks]
\lean{Complexity.BinArith.chunkList}
\label{Complexity.BinArith.chunkList}
\leanok
Given a natural number $c$, a Boolean list $w$, and a natural number $N$, form the list
of length $N$ whose entry at index $i$, for $0\le i<N$, is the length-$c$ segment of $w$
beginning at position $ic$ (or the available portion if fewer than $c$ elements remain).
\end{definition}

\begin{theorem}[Chunk step on encoded pairs]
\lean{Complexity.BinArith.chunkStep_pairEncode}
\label{Complexity.BinArith.chunkStep_pairEncode}
\leanok
\uses{Complexity.BinArith.chunkStep, Complexity.PrefixByLength.drop, Complexity.PrefixByLength.take, Complexity.isNilB, Turing.pairEncode}
\difficulty{1}
For bit strings $c$, $r$, and $a$, apply the chunk step to the encoding of $c$ paired
with the encoding of $r$ paired with $a$. If $r$ is empty, the result is unchanged.
Otherwise, the result pairs $c$ with a pair whose first component is the suffix of $r$
after its first $|c|$ bits and whose second component is $a$ followed by the encoding of
the first $|c|$ bits of $r$.
\end{theorem}

\begin{theorem}[Repeated chunk processing]
\lean{Complexity.BinArith.iterate_chunkStep}
\label{Complexity.BinArith.iterate_chunkStep}
\leanok
\uses{Complexity.BinArith.chunkList, Complexity.BinArith.chunkStep, Complexity.BinArith.chunkStep_pairEncode, Complexity.BinArith.encWords, Complexity.BinArith.encWords_append_singleton, Turing.pairEncode}
\difficulty{6}
Let $c,w$ be bit strings, let $L$ be a list of bit strings, and let $N$ be a natural
number. If the length of $w$ is $N$ times the length of $c$, and either $c$ is nonempty
or $N=0$, then applying the chunk-processing step $N$ times to the pair encoding of $c$
and the pair encoding of $w$ followed by the encoding of $L$ gives the pair encoding of
$c$ and the pair encoding of the empty word followed by the encoding of $L$ with the $N$
consecutive chunks of $w$, each of length equal to the length of $c$, appended.
\end{theorem}

\begin{theorem}[Chunking a paired bit string]
\lean{Complexity.BinArith.chunkW_pairEncode}
\label{Complexity.BinArith.chunkW_pairEncode}
\leanok
\uses{Complexity.BinArith.chunkList, Complexity.BinArith.chunkStep_pairEncode, Complexity.BinArith.chunkW, Complexity.BinArith.encWords, Complexity.BinArith.forLoop_pairEncode, Complexity.BinArith.iterate_chunkStep, Complexity.pairFstD_pairEncode, Complexity.pairSndD_pairEncode, Turing.pairEncode}
\difficulty{4}
Let $c$ and $w$ be finite Boolean words and let $N$ be a natural number. Suppose
$|w|=N|c|$ and either $|c|>0$ or $N=0$. Then applying the chunking operation to the
encoding of $c$ paired with $w$ yields the encoding of the list of $N$ consecutive
chunks of $w$, each of length $|c|$, beginning at positions $0,|c|,\ldots,(N-1)|c|$.
\end{theorem}

\begin{theorem}[Polynomial-time chunk extraction]
\lean{Complexity.BinArith.polyTimeComputable_chunkW}
\label{Complexity.BinArith.polyTimeComputable_chunkW}
\leanok
\uses{Complexity.BinArith.chunkStep, Complexity.BinArith.chunkStep_pairEncode, Complexity.BinArith.chunkW, Complexity.BinArith.forLoop_bound_of_ctx, Complexity.BinArith.pair_cases, Complexity.BinArith.polyTimeComputable_forLoop, Complexity.BinArith.polyTimeComputable_slot2, Complexity.BinArith.polyTimeComputable_slot3, Complexity.PolyTimeComputable, Complexity.PolyTimeComputable.append, Complexity.PolyTimeComputable.comp, Complexity.PolyTimeComputable.pairEncode, Complexity.isNilB, Complexity.pairFstD, Complexity.pairFstD_pairEncode, Complexity.pairSndD, Complexity.pairSndD_pairEncode, Complexity.polyTimeComputable_const, Complexity.polyTimeComputable_dropPrefixByLength, Complexity.polyTimeComputable_id, Complexity.polyTimeComputable_isNil, Complexity.polyTimeComputable_ite, Complexity.polyTimeComputable_pairFstD, Complexity.polyTimeComputable_pairSndD, Complexity.polyTimeComputable_takePrefixByLength, Turing.length_pairEncode, Turing.pairDecode, Turing.pairEncode}
\difficulty{6}
The function obtained by encoding the two projections of an input as nested pairs,
iterating the specified chunk update for one more than the input length, and taking the
second projection three times is polynomial-time computable. In the update, if the first
component of the second projection is empty, the pair is unchanged; otherwise, the first
component is retained and the second component is replaced by the encoding that
rearranges its prefix of that length, the remaining suffix, and the corresponding prefix
of the original second component.
\end{theorem}
